\begin{exercise}{82}
Show that if $f\colon x\to y$ is an equivalence, then the maps
$\mathcal C_{/x}\to\mathcal C_{/y}$ and $\mathcal C_{y/}\to\mathcal C_{x/}$
are Joyal equivalences.

\begin{proof}
Consider the restriction maps
\[
 \mathcal C_{/x}\xleftarrow{a}E=\mathcal C_{/f}
 \xrightarrow{b}\mathcal C_{/y},
 \qquad
 \mathcal C_{y/}\xleftarrow{c}T=\mathcal C_{f/}
 \xrightarrow{d}\mathcal C_{x/}.
\]
Here $a,b$ restrict the final copy of $\Delta^1$ to $x,y$,
and $c,d$ restrict the initial copy to $y,x$, respectively.
All these slices are $\infty$-categories by
Corollary~1.4.24(1).
We claim that all four maps are trivial fibrations.
By the join/slice adjunction (Corollary~1.4.17), lifting a
boundary $\partial\Delta^n\hookrightarrow\Delta^n$ against
each map amounts to the following horn extension in $\mathcal C$:
\[
\begin{array}{c|c|c}
 \text{map}&\text{horn}&\text{reason it fills}\\[2pt]
 a&\Lambda_{n+1}^{n+2}&\text{inner}\\
 b&\Lambda_{n+2}^{n+2}&\text{last edge }f\text{ is an equivalence}\\
 c&\Lambda_1^{n+2}&\text{inner}\\
 d&\Lambda_0^{n+2}&\text{first edge }f\text{ is an equivalence}.
\end{array}
\]
For example, the domain for $a$ is
$(\partial\Delta^n\star\Delta^1)\cup
(\Delta^n\star\{0\})$:
the boundary data prescribe faces $d_0,\ldots,d_n$ and
the specified simplex of $\mathcal C_{/x}$ prescribes $d_{n+2}$.
Replacing $\{0\}$ by $\{1\}$ gives the horn for $b$.
For $c,d$, the boundary data prescribe faces $d_2,\ldots,d_{n+2}$;
restriction to $y,x$ adds $d_0,d_1$, respectively.
The outer horns fill by Theorem~2.1.8 and its opposite form
in Remark~2.1.11, applied to $\mathcal C\to\Delta^0$.
This works for every $n\geq0$; for $n=0$ the horns have
dimension $2$.

By Proposition~2.2.12, all four maps are Joyal equivalences.
Exercise~\ref{ex:59} supplies sections of these trivial
fibrations, since they lift against all monomorphisms.
The induced post-composition functor is $bs$ for any section
$s$ of $a$, as constructed after Corollary~1.4.24:
on an object $u\colon z\to x$, the triangle $s(u)$ has
successive edges $u,f$ and long edge $bs(u)$.
Since $as=\mathrm{id}$ and $a$ is a Joyal equivalence,
$s$ is also a Joyal equivalence, and hence so is $bs$.
Similarly, the induced pre-composition functor is $du$
for any section $u$ of $c$; its triangles have successive
edges $f$ and $y\to z$.
The equality $cu=\mathrm{id}$ makes $u$ a Joyal equivalence,
so $du$ is one as well.
Here the two-out-of-three property is simply that for
isomorphisms in $h\mathrm{Cat}_\infty$
(Definition~2.1.14).
\end{proof}
\end{exercise}
